\section{Force Share by Client Pair}
\label{app:pairs}

\begin{table}[t]
\centering
\caption{Force share of public closes by the client pair at the two ends (closes from 8 May 2023).}
\label{tab:client_pairs}
\footnotesize
\begin{tabular}{lrrr}
\toprule
Client pair & Closes & Force & Force \% \\
\midrule
LND--LND & 169{,}435 & 47{,}905 & 28.3 \\
CLN--LND & 20{,}611 & 7{,}149 & 34.7 \\
LDK--LND & 17{,}749 & 6{,}603 & 37.2 \\
LDK--LDK & 13{,}415 & 204 & 1.5 \\
Eclair--LND & 11{,}474 & 3{,}335 & 29.1 \\
CLN--CLN & 1{,}832 & 740 & 40.4 \\
CLN--Eclair & 927 & 291 & 31.4 \\
\bottomrule
\end{tabular}
\par\smallskip\parbox{\linewidth}{\scriptsize Clients from paper 01's labels on the day of the close (high or medium confidence). The chain does not show which side broadcast a force close, so a close is attributed to the pair.}
\end{table}


Table~\ref{tab:client_pairs} gives the force share of public closes since 8~May~2023 by the implementations at the
two ends of the channel, from the labels of our first paper~\cite{kumar2026clientdiversity}. The chain does not show
which side broadcast a force close, so a close is attributed to the pair, not to one implementation. Among pairs with
both ends labelled and more than 1,000 closes, the force share ranges from \PairForcePctLndLnd\% (LND--LND,
\PairClosesLndLnd{} closes) to \PairForcePctClnCln\% (CLN--CLN); CLN--LND (\PairForcePctClnLnd\%), LDK--LND
(\PairForcePctLdkLnd\%) and Eclair--LND (\PairForcePctEclairLnd\%) lie between. Pairs with an unlabelled end run higher,
for example \PairForcePctLndUnknown\% for LND with an unlabelled peer. The one outlier, LDK--LDK at
\PairForcePctLdkLdkExact\% of \PairClosesLdkLdk{} closes, is one node closing channels cooperatively at scale, not a
property of the implementation: a single node is party to \PairLdkLdkTopNodePct\% of those closes. For the other
pairs we have not separated operators from implementations, so the table describes pairs and does not rank
implementations.
